% ISMRM 2027 abstract draft (Standard abstract)
% Deadline: 2026-10-28 (ECHO platform)
% Limits: body 750 words, up to 5 figures (captions <= 500 characters each),
%         Impact 40 words, Synopsis 100 words (plain text, no equations,
%         references or citations), title 125 characters, plus one Preview
%         Figure (3:2 aspect ratio).
% Compile: pdflatex abstract.tex ; count words: python scripts/count_words.py
%
% Layout follows the order of the ECHO form: title, impact, synopsis, body,
% figure captions, preview figure, references, acknowledgements, keywords.
% \todo{...} marks what must be checked or filled before submission.

\documentclass[11pt]{article}
\usepackage[utf8]{inputenc}
\usepackage[T1]{fontenc}
\usepackage[margin=2cm]{geometry}
\usepackage{graphicx}
\usepackage{xcolor}
\usepackage{booktabs}
\usepackage[hidelinks]{hyperref}

\newcommand{\todo}[1]{\textcolor{red}{[TODO: #1]}}
% Draws a placeholder box until the figure is produced in figures/
\newcommand{\fig}[1]{\IfFileExists{#1}{\includegraphics[width=0.9\linewidth]{#1}}%
  {\fbox{\parbox[c][4cm][c]{0.85\linewidth}{\centering\texttt{\detokenize{#1}} (missing)}}}}

\begin{document}

% ---------------------------------------------------------------------------
\section*{Title (125 characters max)}

% 108 characters
Photogrammetry of brain dissections: automatic mesh closing works, but
acquisition limits volume measurement

% ---------------------------------------------------------------------------
\section*{Impact Statement (40 words max)}

% Word count: scripts/count_words.py
Photogrammetric meshes of brain dissections can be closed automatically, but
their scanned edge moves by millimetres between epochs, changing volumes as
much as one epoch of dissection: acquisition must be controlled before these
meshes serve as quantitative references.

% ---------------------------------------------------------------------------
\section*{Synopsis (100 words max)}

\textbf{Motivation:} Photogrammetric meshes of brain dissections could serve
as quantitative references for diffusion MRI, but they are open surfaces with
scanning defects.

\textbf{Goal(s):} Close them into watertight volumes and measure how far
these volumes can be trusted.

\textbf{Approach:} We repaired and closed all 130 meshes of the public
BraDiPho archive, compared closing methods and checked that volume decreases
along each dissection.

\textbf{Results:} Every mesh became watertight, and closing methods changed
volumes by 0.45\% on average. Yet volume increased in 15 of 104
epoch-to-epoch changes, by up to 15\%, mostly because the scanned edge moved
2 to 5~mm outward.

% ---------------------------------------------------------------------------
\section*{Introduction}

Photogrammetry records a brain dissection as a series of textured 3D
meshes$^{1,2}$: after each removal of tissue, the specimen is scanned again,
giving one mesh per epoch. The BraDiPho archive$^{1}$ shares such meshes to
compare dissected pathways with diffusion MRI tractography, and photographs
of brain slices registered to surface scans already yield regional
volumes$^{3}$. Measuring the volume removed at each epoch requires watertight
surfaces. However, each mesh is open on the unphotographed face and contains scanning defects (inverted faces, folds, floating debris), and a
watertight mesh is not necessarily a correct one (Figure 1).

We closed the whole archive and asked: (1) can the meshes be closed
automatically; (2) how much does the volume depend on how they are closed;
and (3) does the volume decrease at every epoch, as removing tissue
requires? We report what works and what does not.

% ---------------------------------------------------------------------------
\section*{Methods}

\textbf{Data.} We used all 130 epochs of the 12 BraDiPho specimens$^{1}$. Undocumented ``b'' variants were excluded from
comparisons between epochs, leaving 104 epoch-to-epoch changes.

\textbf{Repair and closing} (Python, VTK; Figure 2). (1) Duplicated faces,
small flaps and small fragments are removed, non-manifold edges are split,
and faces are oriented outward. (2) Each opening, bounded by a loop, is
closed by a cap. A plane is fitted to the loop; if the loop projects onto it
without crossing itself (after straightening small folds of up to 20
vertices), a planar cap is triangulated$^{4}$ on the outer side of the
opening. Otherwise, a smooth curved cap is stretched across the loop like a
soap film (harmonic interpolation).
(3) Remaining defects and debris are removed and parts reoriented; volume
is computed only for watertight meshes.

\textbf{Evaluation.} Baselines: a fan closing of the same loops, two
variants of fold straightening (A: always planar; B: planar only if the cap
adds at most 5\% of volume) and MeshFix$^{5}$, an independent repair
tool. We measured (a) watertightness; (b) the closing
uncertainty, i.e.\ the volume difference between methods; (c) visual defects
in every epoch; and (d) plausibility: dissection only removes tissue, so
volume should decrease. Each increase was compared with 7 normal decreases
(controls) for scale (similarity ICP), displacement of
the main cap and volume added by the caps.

% ---------------------------------------------------------------------------
\section*{Results}

\textbf{What works.} All 130 epochs became watertight (Figure 2).
Straightening small folds reduced the epochs needing a curved cap from 38 to
12. Volumes differed from the fan closing by 0.45\% and from MeshFix by
0.34\% on average (maxima 3.4\% and 4.6\%; Figure 3b).
After corrections, the epochs with visual defects dropped from 29 to 12.

\textbf{What works poorly.} Most of these 12 epochs (7 from one specimen)
have truly folded openings, where the curved cap folds onto itself
(Figure 4a); others have small holes just under the planar cap (Figure 4b).
On four deeply hollowed openings, variants A and B differed by 5.7 to 11\%:
there, the volume depends on an arbitrary choice.

\textbf{What does not work.} The median volume removed per epoch was 8.7\%
(range 0.5--37.5\%). Yet volume \emph{increased} in 15 of 104 changes, by up
to 15.3\%, and 21 changes were smaller than the closing uncertainty
(Figure 3a). Scale differences ($\leq$0.8\%) did not explain
them. In 9 of the 15, including the three largest, the main cap lay
2.2 to 5.1~mm further out, adding 19 to 57~cm$^3$, whereas it moved by
$-2.2$ to $+1.9$~mm in the controls (Figure 5). Two others followed a
poorly closed epoch, and four were within the closing uncertainty. MeshFix
found exactly the same 15 increases, up to 15.4\%.

% ---------------------------------------------------------------------------
\section*{Discussion}

Only visual inspection revealed folds and inverted debris. The closing uncertainty (about 3\%, up to 11\%) cannot
explain increases of 13 to 15\%: the scanned edge itself moves, so any
closing method inherits the shift. Each millimetre changes the volume by 2 to
4\%, comparable to the tissue removed in a typical epoch. The limiting factor
is the acquisition, not the closing.

\textbf{With existing data} such as BraDiPho, we suggest planar caps where
possible, a visual check of every epoch, and reporting the closing
uncertainty with each volume, flagging increases and changes within it. \textbf{For new acquisitions},
we recommend photographing both faces at every epoch, fixing the specimen on
a rigid support with markers and a scale reference, and adding
structured-light or lidar scanning, possibly less prone to floating fragments
than stereo photogrammetry. For now, dissection photogrammetry is well suited to
visualization, anatomy teaching and qualitative comparison with tractography;
quantitative volumes are close, but not yet reliable.

% ---------------------------------------------------------------------------
\clearpage
\section*{Figure captions}

% Character counts include spaces (ECHO limit: 500 per caption).

\begin{figure}[h]
  \centering
  \fig{figures/fig1_concept.png}
  \caption*{\textbf{Figure 1} (\texttt{fig1\_concept.png}, 420 characters).
  Why a watertight volume can still be wrong. Cross-section of a dissected
  hemisphere (grey), open on the unphotographed face. (a) The cap (red)
  follows the edge of the scanned surface. (b) If this edge lies 1~mm further
  out at the next epoch, the volume grows by 2--4\%, as much as a typical
  epoch of dissection. (c) A three-vertex fold of the boundary warps a planar
  cap. (d) Inverted debris silently subtracts its volume.}
\end{figure}
\clearpage

\begin{figure}[h]
  \centering
  \fig{figures/fig1_pipeline.png}
  \caption*{\textbf{Figure 2} (\texttt{fig1\_pipeline.png}, 209 characters).
  Pipeline on a successful epoch (spc-01, epoch 1). (a) Textured
  photogrammetric mesh. (b) Open boundary, seen from the unphotographed face
  (inside of the shell in dark grey). (c) Watertight result, caps in red.}
\end{figure}
\clearpage

\begin{figure}[h]
  \centering
  \fig{figures/fig2_volume_vs_epoch.png}
  \caption*{\textbf{Figure 3} (\texttt{fig2\_volume\_vs\_epoch.png}, 482
  characters). (a) Closed volume versus epoch, one panel per specimen. Volume
  should only decrease; red: the 15 increases among the 104 epoch-to-epoch
  changes. Grey band: closing uncertainty, $\pm$3.4\%, the largest difference
  between our closing and a flat fan over all epochs; a change inside it
  cannot be told apart from a closing effect. (b) Volume difference between
  our closing and a flat fan or MeshFix (130 epochs), and the two
  fold-straightening variants (12 problem epochs). Boxes: quartiles.}
\end{figure}
\clearpage

\begin{figure}[h]
  \centering
  \fig{figures/fig3_failures.png}
  \caption*{\textbf{Figure 4} (\texttt{fig3\_failures.png}, 183 characters).
  What works poorly (caps in red, back-facing surfaces in blue). (a) Folded
  opening (spc-07, epoch 2): the curved cap folds. (b) Small holes just under
  the planar cap (spc-15, epoch 8).}
\end{figure}
\clearpage

\begin{figure}[h]
  \centering
  \fig{figures/fig5_scanned_edge.png}
  \caption*{\textbf{Figure 5} (\texttt{fig5\_scanned\_edge.png}, 443
  characters). Where the volume increases come from. (a) Epoch 8 (orange)
  lying outside epoch 7 (grey) of spc-17: the scanned edge moved 4.1~mm
  outward and the volume rose by 15\%. (b) Outward displacement of the main
  cap versus volume change, for all 15 volume increases (orange) and 7 normal
  decreases chosen as controls (grey). Dashed: volume added by the displacement alone (2 and 4\%
  per mm, for cut faces of 70 to 120~cm$^2$); points below also lost
  tissue.}
\end{figure}
\clearpage

% ---------------------------------------------------------------------------
\section*{Preview figure}

\begin{center}
  \fig{figures/preview.png}
\end{center}
\clearpage

% ---------------------------------------------------------------------------
\section*{References (Vancouver, in order of first citation)}

\begin{enumerate}
  \item Vavassori L, Rheault F, Nocerino E, et al. Brain dissection
        photogrammetry: a tool for studying human white matter connections
        integrating ex-vivo and in-vivo multimodal datasets. Nat Commun.
        2025. doi:10.1038/s41467-025-64788-y
  \item De Benedictis A, Nocerino E, Menna F, et al. Photogrammetry of the
        human brain: a novel method for three-dimensional quantitative
        exploration of the structural connectivity in neurosurgery and
        neurosciences. World Neurosurg. 2018;115:e279--e291.
        doi:10.1016/j.wneu.2018.04.036
  \item Gazula H, Tregidgo HFJ, Billot B, et al. Machine learning of
        dissection photographs and surface scanning for quantitative 3D
        neuropathology. eLife. 2024;12:RP91398. doi:10.7554/eLife.91398
  \item Shewchuk JR. Triangle: engineering a 2D quality mesh generator and
        Delaunay triangulator. In: Lin MC, Manocha D, editors. Applied
        Computational Geometry. Lect Notes Comput Sci. 1996;1148:203--222.
        doi:10.1007/BFb0014497
  \item Attene M. A lightweight approach to repairing digitized polygon
        meshes. Vis Comput. 2010;26:1393--1406. doi:10.1007/s00371-010-0416-3
\end{enumerate}

% ---------------------------------------------------------------------------
\section*{Acknowledgements}
% Reviewers do not see this section (ISMRM rules): it is the only place for
% funding, code and data links.
We thank the BraDiPho team for making the dissection photogrammetry data
openly available (https://bradipho.org), and in particular Laurent Petit
(Groupe d'Imagerie Neurofonctionnelle, CNRS, Universit\'e de Bordeaux) for
explaining the acquisition and for helpful discussions on our closing
approaches. The code is available at
https://github.com/guichoquette/close\_mesh. This work was carried out during
an internship funded by the Natural Sciences and Engineering Research Council
of Canada (NSERC), through the research funds of Fran\c{c}ois Rheault,
Universit\'e de Sherbrooke.

\section*{Institutional logo}

\begin{center}
  % Universite de Sherbrooke logo only (SCIL logo goes on the poster);
  % ECHO max 300x300 px
  \includegraphics[height=2.5cm]{figures/logo_udes_300.png}
\end{center}

\section*{Keywords (5)}

Ex-Vivo Applications; Validation; Precision \& Accuracy; Analysis/Processing;
Brain.

% ---------------------------------------------------------------------------
\section*{Authors and categories (ECHO form)}

\begin{itemize}
  \item \textbf{Authors:} Guillaume Choquette (first and presenting author);
        Fran\c{c}ois Rheault. Sherbrooke Connectivity Imaging Laboratory
        (SCIL), Universit\'e de Sherbrooke, Sherbrooke, QC, Canada.
  \item \textbf{Abstract type:} Standard.
  \item \textbf{Primary category:} Analysis Methods -- Data Processing.
  \item \textbf{Secondary category:} Diffusion -- Simulation/Validation.
\end{itemize}

\end{document}
